% ============================================================
%  Programação em Java — Lição 4: Orientação a Objetos — Fundamentos
%  Java Programming MOOC · Universidade de Helsinque
%  Compilar com: pdflatex licao04.tex
% ============================================================
\documentclass[12pt, a4paper]{article}

% ── Codificação e língua ──────────────────────────────────
\usepackage[utf8]{inputenc}
\usepackage[T1]{fontenc}
\usepackage[brazil]{babel}

% ── Fontes ────────────────────────────────────────────────
\usepackage{palatino}
\usepackage[scaled=0.88]{beramono}

% ── Geometria ─────────────────────────────────────────────
\usepackage[
  a4paper,
  top=3cm, bottom=3cm,
  left=3cm, right=3cm
]{geometry}

% ── Cores ─────────────────────────────────────────────────
\usepackage[dvipsnames, table, x11names]{xcolor}

\definecolor{javaOrange}{HTML}{E76F00}
\definecolor{javaDeep}{HTML}{BF5700}
\definecolor{javaBlue}{HTML}{1565C0}
\definecolor{javaTeal}{HTML}{00695C}
\definecolor{javaAmber}{HTML}{B45309}
\definecolor{javaInk}{HTML}{2B2140}
\definecolor{javaCodeBg}{HTML}{FFF8F0}
\definecolor{javaRule}{HTML}{FFD7B5}
\definecolor{javaLightBg}{HTML}{FFF5EC}
\definecolor{javaKeyword}{HTML}{D73A49}
\definecolor{javaString}{HTML}{032F62}
\definecolor{javaComment}{HTML}{6A737D}

% ── Cabeçalho e rodapé ────────────────────────────────────
\usepackage{fancyhdr}
\pagestyle{fancy}
\fancyhf{}
\renewcommand{\headrulewidth}{1.5pt}
\renewcommand{\headrule}{\hbox to\headwidth{\color{javaOrange}\leaders\hrule height \headrulewidth\hfill}}
\fancyhead[L]{\small\color{javaOrange}\textbf{Programação em Java}}
\fancyhead[R]{\small\color{javaDeep}Lição 4: Orientação a Objetos}
\fancyfoot[C]{\small\color{javaAmber}\thepage}
\fancyfoot[R]{\small\color{gray}Java Programming MOOC}

% ── Títulos de seção ──────────────────────────────────────
\usepackage{titlesec}

\titleformat{\section}
  {\color{javaOrange}\large\bfseries}
  {\color{javaOrange}\thesection.}{0.6em}{}
  [\vspace{2pt}\color{javaRule}\titlerule]

\titleformat{\subsection}
  {\color{javaDeep}\normalsize\bfseries}
  {\color{javaDeep}\thesubsection.}{0.5em}{}

\titleformat{\subsubsection}
  {\color{javaTeal}\normalsize\bfseries}
  {\color{javaTeal}\thesubsubsection.}{0.5em}{}

\titlespacing*{\section}{0pt}{2ex}{1ex}
\titlespacing*{\subsection}{0pt}{1.8ex}{0.6ex}
\titlespacing*{\subsubsection}{0pt}{1.4ex}{0.4ex}

% ── Hiperlinks ────────────────────────────────────────────
\usepackage[
  colorlinks=true,
  linkcolor=javaBlue,
  urlcolor=javaBlue,
  citecolor=javaBlue,
  pdfauthor={Java Programming MOOC},
  pdftitle={Licao 4: Orientacao a Objetos -- Fundamentos},
  pdfsubject={Programacao em Java},
  pdfkeywords={java, classes, objetos, encapsulamento, OOP}
]{hyperref}

% ── Código-fonte ─────────────────────────────────────────
\usepackage{listings}

\lstdefinestyle{java}{
  language=Java,
  basicstyle=\small\ttfamily\color{javaInk},
  keywordstyle=\color{javaKeyword}\bfseries,
  stringstyle=\color{javaString},
  commentstyle=\color{javaComment}\itshape,
  numberstyle=\tiny\color{gray},
  numbers=left,
  numbersep=8pt,
  stepnumber=1,
  showstringspaces=false,
  breaklines=true,
  breakatwhitespace=true,
  tabsize=4,
  frame=none,
  morekeywords={var, record, String, Scanner, ArrayList, HashMap,
                System, Integer, Double, Boolean, File, Math,
                Override, FileNotFoundException},
}

\lstdefinestyle{plaintext}{
  basicstyle=\small\ttfamily\color{javaInk},
  numbers=none,
  showstringspaces=false,
  breaklines=true,
  frame=none,
}

% ── tcolorbox para blocos de código ───────────────────────
\usepackage[most]{tcolorbox}
\tcbuselibrary{listings, breakable}

\newtcblisting{javacode}{
  listing only,
  listing style=java,
  breakable,
  enhanced jigsaw,
  colback=javaCodeBg,
  colframe=javaRule,
  borderline west={3pt}{0pt}{javaDeep},
  arc=4pt,
  top=6pt, bottom=6pt, left=10pt, right=6pt,
  boxrule=0.5pt,
}

\newtcblisting{plaincode}{
  listing only,
  listing style=plaintext,
  breakable,
  enhanced jigsaw,
  colback=javaCodeBg,
  colframe=javaRule,
  borderline west={3pt}{0pt}{javaAmber},
  arc=4pt,
  top=6pt, bottom=6pt, left=10pt, right=6pt,
  boxrule=0.5pt,
}

% Caixa de destaque (callout)
\newtcolorbox{callout}[1][Nota]{
  breakable,
  enhanced jigsaw,
  colback=javaLightBg,
  colframe=javaOrange,
  coltitle=white,
  fonttitle=\small\bfseries,
  title=#1,
  arc=4pt,
  boxrule=1pt,
  top=4pt, bottom=4pt, left=8pt, right=8pt,
}

% Caixa de exercício
\newtcolorbox{exercicio}[1]{
  breakable,
  enhanced jigsaw,
  colback=white,
  colframe=javaRule,
  borderline west={4pt}{0pt}{javaOrange},
  coltitle=javaOrange,
  fonttitle=\normalsize\bfseries,
  title=Exercício #1,
  arc=4pt,
  boxrule=0.5pt,
  top=6pt, bottom=6pt, left=10pt, right=8pt,
  attach boxed title to top left={yshift=-2mm, xshift=8pt},
  boxed title style={
    colback=javaOrange, colframe=javaDeep,
    arc=3pt, boxrule=0.5pt,
    fontupper=\small\bfseries\color{white}
  },
}

% ── Tabelas ───────────────────────────────────────────────
\usepackage{booktabs}
\usepackage{longtable}
\usepackage{array}

% ── Utilitários ───────────────────────────────────────────
\usepackage{microtype}
\usepackage{parskip}
\setlength{\parskip}{0.5em}
\usepackage{enumitem}
\setlist[itemize]{topsep=2pt, itemsep=2pt, leftmargin=1.5em}
\setlist[enumerate]{topsep=2pt, itemsep=2pt, leftmargin=1.5em}

% ── Código inline ─────────────────────────────────────────
\newcommand{\code}[1]{\texttt{\color{javaDeep}#1}}

% ──────────────────────────────────────────────────────────
%  CAPA
% ──────────────────────────────────────────────────────────
\usepackage{tikz}

\newcommand{\makejavacapa}{%
  \begin{titlepage}
    \begin{tikzpicture}[remember picture, overlay]
      \fill[javaOrange]
        (current page.north west) rectangle
        ([yshift=-10cm] current page.north east);
      \fill[javaDeep]
        ([yshift=-10cm] current page.north west) rectangle
        ([yshift=-10.5cm] current page.north east);
      \fill[javaAmber!40!white]
        (current page.south west) rectangle
        ([yshift=3cm] current page.south east);
    \end{tikzpicture}

    \vspace*{2.5cm}
    {\color{white}\fontsize{13}{16}\selectfont\textbf{PROGRAMAÇÃO EM JAVA}}\\[0.5cm]
    {\color{white}\fontsize{28}{34}\selectfont\textbf{Lição 4}}\\[0.3cm]
    {\color{white!90!javaOrange}\fontsize{18}{22}\selectfont Orientação a Objetos --- Fundamentos}\\[1.8cm]

    {\color{white!85!javaOrange}\large Java Programming MOOC}\\[0.3cm]
    {\color{white!80!javaOrange}\normalsize Universidade de Helsinque}\\[2cm]

    \vfill
    {\color{javaAmber!60!black}\small
      Tradução para o português do Brasil --- 2026\\
      Licença: CC BY-NC-SA 4.0
    }
    \vspace{1.5cm}
  \end{titlepage}
}

% ══════════════════════════════════════════════════════════
\begin{document}

\makejavacapa
\tableofcontents
\newpage

% ──────────────────────────────────────────────────────────
\section{Objetivos de aprendizagem}
% ──────────────────────────────────────────────────────────

Ao final desta lição, você será capaz de:

\begin{itemize}
  \item Entender o que são classes e objetos
  \item Criar classes com variáveis de instância, construtores e métodos
  \item Usar \code{this} para referenciar o objeto atual
  \item Implementar encapsulamento com \code{private} e métodos getters/setters
  \item Armazenar objetos em listas (\code{ArrayList})
  \item Ler dados de arquivos e representá-los como objetos
\end{itemize}

% ──────────────────────────────────────────────────────────
\section{O que são classes e objetos?}
% ──────────────────────────────────────────────────────────

Uma \textbf{classe} é um molde (blueprint) que descreve um conceito do mundo real: seus atributos e comportamentos. Um \textbf{objeto} é uma instância concreta desse molde.

Analogia: a \textbf{planta baixa} de uma casa é a classe. Cada \textbf{casa construída} a partir dessa planta é um objeto. Todas seguem o mesmo projeto, mas cada uma tem características próprias (cor, moradores, etc.).

\subsection{Por que usar orientação a objetos?}

\begin{itemize}
  \item Modela o problema de forma natural (conceitos do domínio viram classes)
  \item Facilita a organização e manutenção de programas grandes
  \item Promove o reúso de código
\end{itemize}

% ──────────────────────────────────────────────────────────
\section{Criando uma classe}
% ──────────────────────────────────────────────────────────

\begin{javacode}
public class Pessoa {
    // variaveis de instancia (atributos)
    private String nome;
    private int idade;

    // construtor
    public Pessoa(String nome) {
        this.nome = nome;
        this.idade = 0;
    }

    // metodos
    public void crescer() {
        this.idade++;
    }

    public String getNome() {
        return this.nome;
    }

    public int getIdade() {
        return this.idade;
    }

    @Override
    public String toString() {
        return this.nome + ", " + this.idade + " anos";
    }
}
\end{javacode}

\subsection{Usando a classe}

\begin{javacode}
Pessoa alice = new Pessoa("Alice");
Pessoa bruno = new Pessoa("Bruno");

alice.crescer();
alice.crescer();

System.out.println(alice);           // Alice, 2 anos
System.out.println(bruno);           // Bruno, 0 anos
System.out.println(alice.getNome()); // Alice
\end{javacode}

Cada objeto tem seu \textbf{próprio} conjunto de variáveis de instância --- alterar \code{alice} não afeta \code{bruno}.

% ──────────────────────────────────────────────────────────
\section{Variáveis de instância}
% ──────────────────────────────────────────────────────────

As variáveis de instância guardam o \textbf{estado} do objeto. Em geral, são declaradas como \code{private} para protegê-las de acesso externo direto.

\begin{javacode}
public class ContaBancaria {
    private String titular;
    private double saldo;

    public ContaBancaria(String titular, double saldoInicial) {
        this.titular = titular;
        this.saldo = saldoInicial;
    }
}
\end{javacode}

% ──────────────────────────────────────────────────────────
\section{Construtores}
% ──────────────────────────────────────────────────────────

O \textbf{construtor} é chamado com \code{new} e inicializa as variáveis de instância. Tem o mesmo nome da classe e não declara tipo de retorno.

\begin{javacode}
ContaBancaria conta = new ContaBancaria("Maria", 1000.0);
\end{javacode}

\begin{callout}[Importante]
Se você não definir um construtor, Java fornece um padrão (sem parâmetros), mas assim que você define um, o padrão deixa de existir.
\end{callout}

% ──────────────────────────────────────────────────────────
\section{O \texttt{this}}
% ──────────────────────────────────────────────────────────

A palavra \code{this} dentro de um método refere-se ao \textbf{objeto que chamou o método}:

\begin{javacode}
public class Ponto {
    private int x;
    private int y;

    public Ponto(int x, int y) {
        this.x = x; // this.x e a variavel de instancia; x e o parametro
        this.y = y;
    }

    public double distanciaOrigem() {
        return Math.sqrt(this.x * this.x + this.y * this.y);
    }
}
\end{javacode}

% ──────────────────────────────────────────────────────────
\section{Encapsulamento: \texttt{private} e getters/setters}
% ──────────────────────────────────────────────────────────

\textbf{Encapsulamento} significa esconder os detalhes internos do objeto. Declaramos variáveis como \code{private} e fornecemos métodos públicos para lê-las ou alterá-las.

\subsection{Getters (leitura)}

\begin{javacode}
public String getNome() {
    return this.nome;
}

public int getIdade() {
    return this.idade;
}

public boolean isMaiorDeIdade() {
    return this.idade >= 18;
}
\end{javacode}

\subsection{Setters (escrita com validação)}

\begin{javacode}
public void setIdade(int idade) {
    if (idade >= 0) {
        this.idade = idade;
    }
}

public void setSaldo(double saldo) {
    if (saldo >= 0) {
        this.saldo = saldo;
    }
}
\end{javacode}

O setter pode \textbf{validar} o valor antes de atribuí-lo --- isso é impossível com acesso direto à variável.

% ──────────────────────────────────────────────────────────
\section{O método \texttt{toString()}}
% ──────────────────────────────────────────────────────────

O método \code{toString()} define como o objeto é representado como texto. É chamado automaticamente quando o objeto é passado para \code{System.out.println()} ou concatenado com uma string.

\begin{javacode}
@Override
public String toString() {
    return this.titular + " | Saldo: R$ " + this.saldo;
}
\end{javacode}

\begin{javacode}
ContaBancaria conta = new ContaBancaria("Joao", 500.0);
System.out.println(conta); // Joao | Saldo: R$ 500.0
\end{javacode}

% ──────────────────────────────────────────────────────────
\section{Métodos que modificam o estado}
% ──────────────────────────────────────────────────────────

\begin{javacode}
public class ContaBancaria {
    private String titular;
    private double saldo;

    public ContaBancaria(String titular, double saldoInicial) {
        this.titular = titular;
        this.saldo = saldoInicial;
    }

    public void depositar(double valor) {
        if (valor > 0) {
            this.saldo += valor;
        }
    }

    public boolean sacar(double valor) {
        if (valor > 0 && this.saldo >= valor) {
            this.saldo -= valor;
            return true;
        }
        return false;
    }

    public double getSaldo() {
        return this.saldo;
    }

    @Override
    public String toString() {
        return this.titular + " | Saldo: R$ " + String.format("%.2f", this.saldo);
    }
}
\end{javacode}

\begin{javacode}
ContaBancaria conta = new ContaBancaria("Carlos", 200.0);
conta.depositar(300.0);
conta.sacar(100.0);
System.out.println(conta); // Carlos | Saldo: R$ 400,00
\end{javacode}

% ──────────────────────────────────────────────────────────
\section{Objetos em listas}
% ──────────────────────────────────────────────────────────

Objetos podem ser armazenados em \code{ArrayList} como qualquer outro tipo:

\begin{javacode}
ArrayList<Pessoa> pessoas = new ArrayList<>();
pessoas.add(new Pessoa("Ana"));
pessoas.add(new Pessoa("Bob"));
pessoas.add(new Pessoa("Carol"));

for (Pessoa p : pessoas) {
    System.out.println(p);
}
\end{javacode}

\subsection{Buscando um objeto na lista}

\begin{javacode}
public static Pessoa buscarPorNome(ArrayList<Pessoa> lista, String nome) {
    for (Pessoa p : lista) {
        if (p.getNome().equals(nome)) {
            return p;
        }
    }
    return null; // nao encontrado
}
\end{javacode}

% ──────────────────────────────────────────────────────────
\section{Lendo dados de arquivos}
% ──────────────────────────────────────────────────────────

Java pode ler arquivos de texto com \code{Scanner} e \code{File}:

\begin{javacode}
import java.util.Scanner;
import java.io.File;
import java.io.FileNotFoundException;

public class LeitorDeArquivo {
    public static void main(String[] args) throws FileNotFoundException {
        Scanner scanner = new Scanner(new File("pessoas.txt"));

        ArrayList<Pessoa> pessoas = new ArrayList<>();

        while (scanner.hasNextLine()) {
            String linha = scanner.nextLine();
            String[] partes = linha.split(",");
            String nome = partes[0];

            Pessoa p = new Pessoa(nome);
            pessoas.add(p);
        }

        for (Pessoa p : pessoas) {
            System.out.println(p);
        }
    }
}
\end{javacode}

Exemplo de arquivo \code{pessoas.txt}:

\begin{plaincode}
Alice,30
Bruno,25
Carol,28
\end{plaincode}

% ──────────────────────────────────────────────────────────
\section{Exemplo completo: classe \texttt{Produto}}
% ──────────────────────────────────────────────────────────

\begin{javacode}
public class Produto {
    private String nome;
    private double preco;
    private int estoque;

    public Produto(String nome, double preco, int estoque) {
        this.nome = nome;
        this.preco = preco;
        this.estoque = estoque;
    }

    public String getNome() { return this.nome; }
    public double getPreco() { return this.preco; }
    public int getEstoque() { return this.estoque; }

    public boolean vender(int quantidade) {
        if (quantidade > this.estoque) {
            System.out.println("Estoque insuficiente!");
            return false;
        }
        this.estoque -= quantidade;
        return true;
    }

    public void repor(int quantidade) {
        this.estoque += quantidade;
    }

    @Override
    public String toString() {
        return this.nome + " | R$ " + String.format("%.2f", this.preco)
               + " | Estoque: " + this.estoque;
    }
}
\end{javacode}

% ──────────────────────────────────────────────────────────
\section{Resumo}
% ──────────────────────────────────────────────────────────

\begin{itemize}
  \item Uma \textbf{classe} define o molde; um \textbf{objeto} é uma instância da classe.
  \item \textbf{Variáveis de instância} (\code{private}) guardam o estado do objeto.
  \item O \textbf{construtor} inicializa o objeto ao criá-lo com \code{new}.
  \item \textbf{\code{this}} refere-se ao próprio objeto dentro dos seus métodos.
  \item \textbf{Encapsulamento}: atributos \code{private}, acesso via getters e setters.
  \item \textbf{\code{toString()}} define a representação textual do objeto.
  \item Objetos podem ser armazenados em \code{ArrayList<NomeDaClasse>}.
\end{itemize}

% ──────────────────────────────────────────────────────────
\section{Exercícios}
% ──────────────────────────────────────────────────────────

\begin{exercicio}{1 --- Retângulo}
Crie a classe \code{Retangulo} com atributos \code{largura} e \code{altura} (ambos \code{double}, \code{private}). Adicione:
\begin{itemize}
  \item Construtor com os dois valores
  \item Métodos \code{getArea()} e \code{getPerimetro()}
  \item Método \code{ehQuadrado()} que retorna \code{boolean}
  \item \code{toString()} descritivo
\end{itemize}
\end{exercicio}

\begin{exercicio}{2 --- Aluno e notas}
Crie a classe \code{Aluno} com \code{nome} (String) e uma lista de notas (\code{ArrayList<Integer>}). Adicione:
\begin{itemize}
  \item Método \code{adicionarNota(int nota)} (valide: nota entre 0 e 10)
  \item Método \code{calcularMedia()} retornando \code{double}
  \item Método \code{aprovado()} que retorna \code{true} se média $\geq$ 6
  \item \code{toString()} mostrando nome, média e situação
\end{itemize}
\end{exercicio}

\begin{exercicio}{3 --- Estacionamento}
Crie a classe \code{Veiculo} com \code{placa} (String) e \code{marca} (String). Crie a classe \code{Estacionamento} com:
\begin{itemize}
  \item Capacidade máxima (definida no construtor)
  \item Lista de veículos
  \item Métodos \code{entrar(Veiculo v)} (verifica se há vaga), \code{sair(String placa)}, \code{listarVeiculos()}
\end{itemize}
\end{exercicio}

\begin{exercicio}{4 --- Leitura de arquivo}
Crie um arquivo \code{produtos.txt} com linhas no formato \code{nome,preco,estoque}.
Leia o arquivo, crie objetos \code{Produto} e exiba todos os produtos. Calcule o valor total em estoque (preço $\times$ quantidade).
\end{exercicio}

% ──────────────────────────────────────────────────────────
\vfill
\begin{center}
  \small\color{gray}
  Java Programming MOOC · Universidade de Helsinque ·
  \href{https://java-programming.mooc.fi/part-4}{java-programming.mooc.fi/part-4}\\
  Licença: \href{https://creativecommons.org/licenses/by-nc-sa/4.0/}{CC BY-NC-SA 4.0} ·
  Tradução para o português do Brasil --- 2026
\end{center}

\end{document}
